\section{Experiments}
\label{sec:experiments}

\subsection{Experiment Setup}
\label{sec:experiment_setup}

\subsubsection{Datasets}
\label{sec:datasets}

We evaluate \method{} on two publicly available 3D chest CT report generation benchmarks.
\textbf{CT-RATE}~\cite{hamamci2024ct2rep} pairs 3D chest CT volumes with free-text radiology reports.
\textbf{RadGenome-ChestCT}~\cite{zhang2024radgenome} provides 3D chest CT volumes with structured radiology reports that contain finer-grained anatomical annotations.
From each dataset, we sample 2{,}500 cases and split them into train/valid/test sets with an 80/10/10 ratio (seed$=$42), yielding 2{,}000/250/250 per source.
Four corrupted volumes are excluded, leaving 496 test cases (3{,}979 sentences) for the final evaluation.
All volumes are loaded with SimpleITK, intensity-normalized to $[0,1]$, and trilinearly resized to $128\times128\times128$.

\subsubsection{Baselines}
\label{sec:baselines}

We compare \method{} against four published 3D CT report generation methods:
\textbf{CT2Rep}~\cite{hamamci2024ct2rep}, which decodes reports from a 3D visual prefix;
\textbf{CT-AGRG}~\cite{dipiazza2024ctagrg}, which augments CT2Rep with anatomical graph reasoning;
\textbf{3D-CT-GPT}~\cite{chen2024_3dctgpt}, which adapts GPT-style autoregressive decoding to 3D CT;
and \textbf{Reg2RG}~\cite{chen2025reg2rg}, which introduces region-level grounding for report generation on RadGenome.
For RadGenome, we additionally include \textbf{MedVInT}~\cite{zhang2024medvint} as reported by the Reg2RG authors.
As a provenance-interface reference from the 2D domain, we also include \textbf{MAIRA-2}~\cite{bannur2024maira2}, a grounded CXR report generation system.
None of the 3D baselines expose sentence-conditioned cited evidence objects in their output; rule-based spatial verification and repair are therefore not applicable to them.

\subsubsection{Implementation Details}
\label{sec:implementation}

The 3D encoder is a frozen SwinUNETR~\cite{hatamizadeh2022swinunetr} (62M parameters, self-supervised pretraining on BTCV~\cite{tang2022selfsupervised}) from MONAI~\cite{cardoso2022monai}, producing 48-dimensional features.
The text encoder is a frozen all-MiniLM-L6-v2~\cite{reimers2019sentencebert,wang2020minilm} (22M parameters, 384-dimensional embeddings).
The only trainable component is the cross-modal projection $\mathbf{W}_{\mathrm{proj}}\in\mathbb{R}^{384\times 768}$ (295K parameters), optimized with an InfoNCE objective~\cite{oord2018cpc} using AdamW ($\mathrm{lr}{=}10^{-3}$, cosine schedule, up to 50 epochs, early stopping on validation loss).
For sentence generation and LLM-based judging, we use Llama-3.1-8B-Instruct~\cite{touvron2023llama} in bfloat16 with temperature$=$0.3 and max tokens$=$256.
All configurations share the following hyperparameters: token bank size $B{=}128$, per-sentence citation budget $k{=}8$, spatial IoU weight $\lambda_{\mathrm{spatial}}{=}0.3$, boundary blending $\beta{=}0.1$, and input volume size $128^3$.

We report three layers of evaluation metrics.
\emph{Spatial consistency}: the sentence-level violation rate (Viol\%), defined as the fraction of sentences with at least one \textsc{CLEAR}-6 violation, along with per-rule counts (R1, R3, R6b).
\emph{NLG quality}: BLEU-4~\cite{papineni2002bleu}, ROUGE-L~\cite{lin2004rouge}, and METEOR~\cite{banerjee2005meteor} computed at the sentence level.
\emph{Grounding metrics}: mean anatomy overlap, Hit@$k$, routing accuracy, laterality accuracy, citation faithfulness (CF), and depth faithfulness (DF).
Statistical significance is assessed via patient-level paired bootstrap (5{,}000 resamples, 95\% CI) with Holm step-down correction for multiple comparisons.


\subsection{Performance Comparison}
\label{sec:performance}

Table~\ref{tab:comparison} compares \method{} with published baselines on both datasets.
We classify each method by its output-level provenance interface: whether the system exposes per-sentence cited evidence objects, whether evidence is mandatory or optional, and whether deterministic verification and closed-loop repair are applicable.
All existing 3D CT generators use dense visual features internally but do not expose sentence-conditioned cited 3D evidence objects in their output interface.
\method{} (B2\textquotesingle v2) achieves competitive NLG scores on both datasets while being the only method that provides mandatory token-level citations with executable verification rules.
We note that the NLG numbers are obtained under sentence-prompted generation (where the planner uses reference sentences as topics), and hence constitute controlled upper bounds rather than free-form generation scores.
The evaluation protocols also differ between baselines (published splits) and \method{} (our train/valid/test split), so direct NLG comparison requires caution.

% Comparison table for ProveTok paper
% Requires: \usepackage{booktabs, amssymb, pifont, multirow}
% Define in preamble:
%   \newcommand{\cmark}{\ding{51}}
%   \newcommand{\xmark}{\ding{55}}

\begin{table*}[t]
\centering
\caption{Comparison with existing methods for 3D CT report generation. Citation granularity, evidence
enforcement, verification, and repair-loop columns characterise provenance capabilities.
NLG metrics are BLEU-4, ROUGE-L, and METEOR (higher is better).}
\label{tab:comparison}
\small
\setlength{\tabcolsep}{4pt}
\begin{tabular}{@{}lccccccccc@{}}
\toprule
\textbf{Method} & \textbf{Input / Data} & \textbf{Citation Gran.} & \textbf{Evid.} & \textbf{Verification} & \textbf{Repair} & \textbf{B-4} & \textbf{R-L} & \textbf{MTR} \\
\midrule
% ---- Block A: 3D CT report generation baselines ----
CT2Rep~\cite{hamamci2024ct2rep}           & 3D CT / CT-RATE   & none            & \xmark & none             & \xmark   & 0.369          & 0.459          & 0.295 \\
CT-AGRG~\cite{dipiazza2024ctagrg}         & 3D CT / CT-RATE   & none            & \xmark & none             & \xmark   & 0.172          & 0.280          & 0.196 \\
3D-CT-GPT~\cite{chen2024_3dctgpt}         & 3D CT / CT-RATE   & none            & \xmark & none             & \xmark   & 0.384\textsuperscript{*} & 0.328 & 0.357 \\
MedVInT~\cite{zhang2024medvint}           & 3D CT / RadGenome  & none            & \xmark & none             & \xmark   & 0.246          & 0.326          & 0.404 \\
CT2Rep~\cite{hamamci2024ct2rep}           & 3D CT / RadGenome  & none            & \xmark & none             & \xmark   & 0.236          & 0.310          & 0.402 \\
Reg2RG~\cite{chen2025reg2rg}              & 3D CT / RadGenome  & partial\textsuperscript{\S} & \xmark & none & \xmark   & 0.249          & 0.367          & 0.441 \\
\midrule
% ---- Block B: Grounded reference system (2D) ----
MAIRA-2\textsuperscript{\dag}~\cite{bannur2024maira2} & 2D CXR / MIMIC-CXR & sentence-level  & \xmark & post-hoc scorer  & \xmark   & ---            & ---            & --- \\
\midrule
% ---- Block C: Ours ----
\textbf{ProveTok (B2'v2)} & \textbf{3D CT / CT-RATE}  & \textbf{token citation} & \textbf{\cmark} & \textbf{R1--R6 rules} & \textbf{optional} & \textbf{0.467\textsuperscript{\ddag}} & \textbf{0.626\textsuperscript{\ddag}} & \textbf{0.603\textsuperscript{\ddag}} \\
\textbf{ProveTok (B2'v2)} & \textbf{3D CT / RadGenome} & \textbf{token citation} & \textbf{\cmark} & \textbf{R1--R6 rules} & \textbf{optional} & \textbf{0.506\textsuperscript{\ddag}} & \textbf{0.680\textsuperscript{\ddag}} & \textbf{0.660\textsuperscript{\ddag}} \\
\bottomrule
\end{tabular}

\vspace{4pt}
{\footnotesize\raggedright
\textsuperscript{*}3D-CT-GPT reports BLEU (not BLEU-4); not directly comparable.\\
\textsuperscript{\dag}MAIRA-2 is a 2D CXR system; included as capability reference only.\\
\textsuperscript{\ddag}ProveTok operates in a sentence-prompted setting; NLG scores are controlled upper bounds.\\
\textsuperscript{\S}``partial'': region-level grounding without token-level sentence citation.\\
RadGenome baseline numbers (MedVInT, CT2Rep, Reg2RG) reproduced by Reg2RG authors~\cite{chen2025reg2rg}.\par
}
\end{table*}



\subsubsection{Ablation Study}
\label{sec:ablation}

Table~\ref{tab:ablation} presents the seven-configuration ablation chain evaluated on the full 5K test set (3{,}979 sentences from 496 patients).
Each configuration adds exactly one component to the previous one, isolating the marginal contribution of each module.

The chain reveals two statistically significant drops in violation rate.
First, training $\mathbf{W}_{\mathrm{proj}}$ (A0$\to$A1) reduces Viol\% from 7.59\% to 6.01\% ($p_{\mathrm{adj}}{<}0.01$), as the learned projection better aligns visual tokens with sentence queries and reduces depth mismatches (R3: 237$\to$174).
Second, switching from Evidence Card~v1 to Evidence Card~v2 (B2\textquotesingle$\to$B2\textquotesingle v2) yields the largest single-step improvement, with Viol\% dropping from 6.01\% to 4.25\% ($p_{\mathrm{adj}}{=}0.006$).
The strict laterality and depth gates of Evidence Card~v2 directly prevent the LLM from generating spatially inconsistent claims, with R3 falling from 86 to 23.

The spatial-filter semantic-rerank routing (A1$\to$E1) shifts the violation profile: R3 drops from 174 to 86 but R1 rises from 0 to 95, reflecting a trade-off between depth and laterality consistency.
Post-hoc LLM judging (C2\textquotesingle) and repair (D2), branching from B2\textquotesingle{} with Evidence Card~v1, do not improve over B2\textquotesingle v2, confirming that constructive evidence gating outperforms post-hoc correction.
NLG metrics remain stable across all generation configurations (BLEU-4 varies by $<$0.01), indicating that evidence gating does not degrade text quality.

% Ablation table for ProveTok paper
% Requires: \usepackage{booktabs, pifont}
% Define in preamble:
%   \newcommand{\cmark}{\ding{51}}
%   \newcommand{\xmark}{\ding{55}}

\begin{table*}[t]
\centering
\caption{Ablation study on the 5K test set (3{,}979 sentences from 496 patients).
Viol\% is the sentence-level spatial violation rate (lower is better).
95\% CI is obtained via patient-level paired bootstrap ($R{=}5{,}000$).
R1/R3/R6b are per-rule violation counts.
NLG metrics are reported only for configurations with LLM generation.}
\label{tab:ablation}
\small
\setlength{\tabcolsep}{4pt}
\begin{NoHyper}
\begin{tabular}{@{}clcccccccc@{}}
\toprule
& \textbf{Configuration} & \textbf{Viol\%\,$\downarrow$} & \textbf{95\% CI} & \textbf{R1} & \textbf{R3} & \textbf{R6b} & \textbf{B-4} & \textbf{R-L} & \textbf{MTR} \\
\midrule
\multicolumn{10}{@{}l}{\textit{Routing only (no LLM generation)}} \\[2pt]
A0 & Identity $W$ + anatomy routing           & 7.59 & [6.84, 8.35] &   0 & 237 & 65 & {--} & {--} & {--} \\
A1 & Trained $W_{\mathrm{proj}}$ + anatomy routing & 6.01\textsuperscript{*} & [5.27, 6.64] &   0 & 174 & 65 & {--} & {--} & {--} \\
E1 & Spatial filter + semantic rerank          & 6.18 & [5.34, 6.77] &  95 &  86 & 65 & {--} & {--} & {--} \\
\midrule
\multicolumn{10}{@{}l}{\textit{+ LLM generation with evidence card}} \\[2pt]
B2\textquotesingle    & \quad + evidence card v1                  & 6.01 & [5.14, 6.60] & 109 &  86 & 44 & 0.471 & 0.643 & 0.618 \\
\textbf{B2\textquotesingle v2} & \quad\textbf{+ evidence card v2 (strict)}  & \textbf{4.25}\textsuperscript{*} & \textbf{[3.60, 4.83]} & \textbf{98} & \textbf{23} & \textbf{48} & \textbf{0.473} & \textbf{0.643} & \textbf{0.619} \\
\midrule
\multicolumn{10}{@{}l}{\textit{+ Post-hoc verification (from B2\textquotesingle)}} \\[2pt]
C2\textquotesingle    & \quad + LLM judge                         & 6.18\textsuperscript{\dag} & [5.25, 6.79] & 113 &  86 & 47 & 0.477 & 0.643 & 0.622 \\
D2   & \quad\quad + repair executor              & 5.96 & [5.07, 6.55] & 102 &  86 & 49 & 0.474 & 0.643 & 0.622 \\
\bottomrule
\end{tabular}
\end{NoHyper}

\vspace{4pt}
{\footnotesize\raggedright
\textsuperscript{*}Statistically significant improvement over the preceding row ($p_{\mathrm{adj}}{<}0.01$, Holm-corrected permutation test, $10{,}000$ permutations).\\
\textsuperscript{\dag}Significant \emph{regression} vs.\ B2\textquotesingle v2 ($p_{\mathrm{adj}}{=}0.006$): the LLM judge occasionally re-routes tokens into new violations.\\
B2\textquotesingle v2 branches from B2\textquotesingle{} with a stricter evidence card; C2\textquotesingle/D2 branch from B2\textquotesingle{} with post-hoc verification.
NLG scores are controlled upper bounds under sentence-prompted generation.\par
}
\end{table*}



\subsubsection{Spatial Grounding and Generalization}
\label{sec:grounding}

Table~\ref{tab:grounding} evaluates spatial grounding quality and citation faithfulness across configurations.
The anatomy-primary routing baselines (A0, A1) achieve perfect overlap and Hit@$k$ by construction, since IoU directly controls token selection.
The semantic-rerank configurations (E1 and beyond) trade overlap for semantic relevance: mean overlap decreases from 0.931 to 0.685, but Hit@8 remains above 98\%, indicating that at least one well-localized token is consistently routed.

Evidence Card~v2 (B2\textquotesingle v2) improves both citation faithfulness (CF: 76.9\%$\to$77.6\%) and depth faithfulness (DF: 97.3\%$\to$98.4\%) over v1, confirming that the strict cleanup step removes misleading tokens and produces more faithful generations.
The post-hoc judge (C2\textquotesingle) achieves the highest CF (81.1\%), as the LLM judge selectively confirms and repairs laterality mismatches; however, this comes at the cost of additional LLM calls and a regression in Viol\% (Table~\ref{tab:ablation}).

% Grounding & Citation Faithfulness table for ProveTok paper
% Requires: \usepackage{booktabs}

\begin{table*}[t]
\centering
\caption{Spatial grounding and citation faithfulness across ablation configurations,
evaluated on 564 grounding-eligible sentences (bilateral cases excluded).
Overlap is the mean intersection-over-token-volume between cited tokens and the target anatomy.
Hit@$k$ measures whether at least one of the top-$k$ cited tokens overlaps ${\geq}50\%$ with the anatomy region.
CF (Citation Faithfulness) and DF (Depth Faithfulness) are only available for LLM-generated configurations.}
\label{tab:grounding}
\small
\setlength{\tabcolsep}{4.5pt}
\begin{NoHyper}
\begin{tabular}{@{}clcccccccc@{}}
\toprule
& \textbf{Configuration} & \textbf{Overlap} & \textbf{Hit@1} & \textbf{Hit@4} & \textbf{Hit@8} & \textbf{Routing Acc.} & \textbf{Lat.\ Acc.} & \textbf{CF} & \textbf{DF} \\
\midrule
\multicolumn{10}{@{}l}{\textit{Routing only (no LLM generation)}} \\[2pt]
A0 & Identity $W$ + anatomy routing      & .931 & 100.0 & 100.0 & 100.0 & 100.0 & 100.0 & {--} & {--} \\
A1 & Trained $W_{\mathrm{proj}}$ + anatomy routing & .931 & 100.0 & 100.0 & 100.0 & 100.0 & 100.0 & {--} & {--} \\
E1 & Spatial filter + semantic rerank     & .685 &  76.8 &  96.1 &  99.3 &  75.4 &  97.2 & {--} & {--} \\
\midrule
\multicolumn{10}{@{}l}{\textit{+ LLM generation with evidence card}} \\[2pt]
B2\textquotesingle    & \quad + evidence card v1              & .685 &  76.8 &  96.1 &  99.3 &  75.4 &  96.9 & 76.9 & 97.3 \\
\textbf{B2\textquotesingle v2} & \quad\textbf{+ evidence card v2 (strict)} & \textbf{.689} & \textbf{76.8} & \textbf{95.9} & \textbf{98.2} & \textbf{75.8} & \textbf{97.1} & \textbf{77.6} & \textbf{98.4} \\
\midrule
\multicolumn{10}{@{}l}{\textit{+ Post-hoc verification (from B2\textquotesingle)}} \\[2pt]
C2\textquotesingle    & \quad + LLM judge                    & .685 &  76.8 &  96.1 &  99.3 &  75.4 &  96.7 & 81.1 & 97.3 \\
D2   & \quad\quad + repair executor         & .685 &  76.8 &  96.1 &  99.3 &  75.4 &  97.0 & 80.6 & 97.3 \\
\bottomrule
\end{tabular}
\end{NoHyper}

\vspace{4pt}
{\footnotesize\raggedright
\textbf{Routing Acc.}: fraction of cited tokens falling within the correct anatomy region.
\textbf{Lat.\ Acc.}: fraction of lateralised sentences free of R1 violations.
\textbf{CF}: generated text laterality agrees with the cited token positions.
\textbf{DF}: token depth falls within the expected anatomical range ($n{=}1{,}423$).\\
A0/A1 achieve perfect spatial metrics by construction (anatomy-priority routing returns only in-region tokens);
the E1+ configurations trade overlap for semantic relevance via cross-modal reranking.\par
}
\end{table*}



\subsubsection{Qualitative Analysis}
\label{sec:qualitative}

Figure~\ref{fig:case_study} presents a representative closed-loop audit trace from the test set.
In this example, the initial generation (Stage~3c) produces a sentence claiming ``left lower lobe'' findings, but \textsc{CLEAR}-6 (Stage~4) detects an R1 laterality violation: the cited tokens are predominantly located on the right side (same-side ratio $=$ 0.25).
The LLM judge (Stage~5a) confirms the violation and suggests a \texttt{reroute\_same\_side} action.
The repair executor (Stage~5b) filters the citation set to right-side tokens only, rebuilds the Evidence Card, and regenerates the sentence as ``right lower lobe,'' which passes all verifier rules.
This example illustrates how the constructive citation interface enables deterministic error detection and targeted repair without modifying the underlying LLM.

\begin{figure}[t]
\centering
% \includegraphics[width=\linewidth]{figures/case_study.pdf}
\caption{Closed-loop audit trace for a laterality violation. Left: initial generation with R1 violation (cited tokens on wrong side). Center: LLM judge confirms the violation. Right: after same-side rerouting and regeneration, the sentence passes all \textsc{CLEAR}-6 rules.}
\label{fig:case_study}
\end{figure}
